\chapter{数列初步}\label{ch:g11:seq}

\omterm{def:g11:seq:sequence}{数列}是由某个规则产生的一列数：例如储蓄账户逐期的余额，
或某种群逐年的规模。本章研究应用中最重要的两类
\omterm{def:g11:seq:sequence}{数列}——
\emph{\omterm{def:g11:seq:arithmetic}{等差数列}}，它们按相等的步长增长，以及
\emph{\omterm{def:g11:seq:geometric}{等比数列}}，它们按相等的比值增长。严格的极限理论
在\cref{ch:g12:seq}中建立。

\section{数列的定义}

\begin{definition}[数列]\label{def:g11:seq:sequence}
称 $(u_n)$ 为\emph{数列}\index{数列}，它对每个\omterm{def:g10:numbers:sets}{整数}
$n \geq 0$（或 $n \geq 1$）对应一个\omterm{def:g10:numbers:sets}{实数} $u_n$，称其为\emph{第
$n$ 项}。数列可以由以下两种方式给出：
\begin{itemize}
  \item \emph{显式地}，由 $u_n$ 关于 $n$ 的公式给出，例如\
        $u_n = n^2 + 1$；
  \item \emph{递推地}，由首项以及由每一项得到下一项的规则给出，
        例如 $u_0 = 3$，$u_{n+1} = 2u_n - 1$。
\end{itemize}
\end{definition}

\begin{example}
对 $u_n = n^2 + 1$，有 $u_0 = 1$，$u_1 = 2$，$u_2 = 5$，$u_{10} = 101$
可直接算出。对 $u_0 = 3$，$u_{n+1} = 2u_n - 1$，有 $u_1 = 5$，$u_2 = 9$，
$u_3 = 17$——每一项都需要前一项；求 $u_{10}$ 需要
十步递推（或用通项公式，见\cref{exo:g11:seq:11}）。
\end{example}

\section{等差数列}

\begin{definition}[等差数列]\label{def:g11:seq:arithmetic}
称\omterm{def:g11:seq:sequence}{数列}为\emph{等差数列}\index{等差数列}，
若其\emph{公差}\index{公差}为 $d$，即每一项都由前一项
加上 $d$ 得到：
\[
u_{n+1} = u_n + d \quad \text{ 对所有 } n.
\]
等价地：差 $u_{n+1} - u_n$ 为常数，且等于 $d$。
\end{definition}

\begin{theorem}[通项公式]\label{thm:g11:seq:arithgeneral}
若 $(u_n)$ 是\omterm{def:g11:seq:arithmetic}{等差数列}，首项为 $u_0$，\omterm{def:g11:seq:arithmetic}{公差}为 $d$，
则
\[
u_n = u_0 + n\,d \quad \text{ 对所有 } n \geq 0,
\qquad\text{ 且更一般地 } u_n = u_p + (n - p)\,d .
\]
\end{theorem}

\begin{proof}
由 $u_0$ 到 $u_n$，共应用“加上 $d$”这一规则 $n$ 次，一
步得 $u_1 = u_0 + d$，两步得 $u_2 = u_0 + 2d$，经过 $n$
步后，每次应用都贡献了一个 $d$，因此 $u_n = u_0 + nd$。
（这一“依此类推”的说法，在\cref{ch:g12:seq}中借助归纳法得到严格化。）
一般公式可通过计算 $n - p$ 步（即从 $u_p$ 到
$u_n$ 的步数）得到。
\end{proof}

\begin{theorem}[连续整数之和]\label{thm:g11:seq:intsum}
对每个\omterm{def:g10:numbers:sets}{整数} $n \geq 1$，有：
\[
1 + 2 + \dots + n = \frac{n(n+1)}{2}.
\]
更一般地，\omterm{def:g11:seq:arithmetic}{等差数列}中连续若干项的和
等于
\[
(\text{项数}) \times
\frac{\text{首项} + \text{末项}}{2}.
\]
\end{theorem}

\begin{proof}
将和 $S$ 写两次，第二次按相反顺序书写，并
逐列相加：
\[
\begin{array}{ccccccccc}
S & = & 1 & + & 2 & + & \dots & + & n\\
S & = & n & + & (n-1) & + & \dots & + & 1\\
\hline
2S & = & (n+1) & + & (n+1) & + & \dots & + & (n+1)
\end{array}
\]
共有 $n$ 列，每列之和为 $n + 1$，故 $2S = n(n+1)$。对于
一般的\omterm{def:g11:seq:arithmetic}{等差数列}，同样的配对依然成立：首项 $+$ 末项 $=$
第二项 $+$ 倒数第二项 $= \dots$，因为左端前进一步
（$+d$）被右端后退一步（
$-d$）所抵消。
\end{proof}

\begin{example}
$1 + 2 + \dots + 100 = \frac{100 \times 101}{2} = 5050$。奇数
之和 $1 + 3 + \dots + 99$（共 $50$ 项）为
$50 \times \frac{1 + 99}{2} = 2500$。
\end{example}

\section{等比数列}

\begin{definition}[等比数列]\label{def:g11:seq:geometric}
称\omterm{def:g11:seq:sequence}{数列}为\emph{等比数列}\index{等比数列}，若其
\emph{公比}\index{公比} $q \neq 0$，即每一项都由前一项
乘以 $q$ 得到：
\[
u_{n+1} = q\,u_n \quad \text{ 对所有 } n.
\]
等价地，当各项均不为零时：比 $\frac{u_{n+1}}{u_n}$ 为
常数，且等于 $q$。
\end{definition}

\begin{theorem}[通项公式]\label{thm:g11:seq:geomgeneral}
若 $(u_n)$ 为\omterm{def:g11:seq:geometric}{等比数列}，首项为 $u_0$，\omterm{def:g11:seq:geometric}{公比}为 $q$，则
\[
u_n = u_0\, q^n \quad \text{对所有 } n \geq 0,
\qquad\text{更一般地 } u_n = u_p\, q^{\,n-p} .
\]
\end{theorem}

\begin{proof}
与\cref{thm:g11:seq:arithgeneral}中的计数方法相同：从 $u_0$ 到
$u_n$，规则“乘以 $q$”共应用 $n$ 次，
贡献因子 $q^n$。
\end{proof}

\begin{theorem}[等比数列求和]\label{thm:g11:seq:geomsum}
对任意\omterm{def:g10:numbers:sets}{实数} $q \neq 1$ 及\omterm{def:g10:numbers:sets}{整数} $n \geq 0$：
\[
1 + q + q^2 + \dots + q^n = \frac{1 - q^{\,n+1}}{1 - q}.
\]
\end{theorem}

\begin{proof}
设 $S = 1 + q + \dots + q^n$，两边乘以 $q$：
$qS = q + q^2 + \dots + q^{n+1}$。两式相减：
\[
S - qS = \bigl(1 + q + \dots + q^n\bigr)
- \bigl(q + q^2 + \dots + q^{n+1}\bigr) = 1 - q^{\,n+1},
\]
因为两式中每个中间项各出现一次，并相互抵消。
于是 $(1 - q)S = 1 - q^{\,n+1}$，两边除以 $1 - q \neq 0$，得到
该公式。
\end{proof}

\begin{example}\label{ex:g11:seq:chessboard}
$1 + 2 + 4 + \dots + 2^{10} = \frac{1 - 2^{11}}{1 - 2} = 2^{11} - 1 =
2047$：若在棋盘的各格上逐格将稻粒数加倍，
则远在第 $64$ 格之前，稻粒数便会超过任何粮仓的容量，
而到该格时总数为 $2^{64} - 1 \approx 1.8 \times 10^{19}$。
\end{example}

\begin{omfigure}
\begin{tikzpicture}
\begin{axis}[omaxis,
  width=10cm, height=5.6cm,
  xlabel={$n$}, ylabel={$u_n$},
  xmin=-0.5, xmax=10.8, ymin=0, ymax=13,
  xtick={5,10}, ytick={2,6}]
  \addplot[omDef, only marks, mark=*, mark size=1.5pt,
    samples at={0,...,10}] {2 + 0.9*x};
  \addplot[omThm, only marks, mark=*, mark size=1.5pt,
    samples at={0,...,10}] {2*1.2^x};
  \node[omDef, font=\small, anchor=east] at (axis cs:9.7,10.9)
    {等差数列};
  \node[omThm, font=\small, anchor=west] at (axis cs:8.6,7.1)
    {等比数列};
\end{axis}
\end{tikzpicture}

{\small 等差与等比的对比：\omterm{def:g11:seq:arithmetic}{等差数列}
（$u_{n+1} = u_n + 0.9$，蓝色）沿直线变化，\omterm{def:g11:seq:geometric}{等比数列}
（$u_{n+1} = 1.2\,u_n$，红色）沿指数曲线变化，
最终超过前者。}
\end{omfigure}

\begin{method}[判断数列的类型]\label{met:g11:seq:recognize}
计算 $u_{n+1} - u_n$ 并化简。若结果为常数 $d$，则
\omterm{def:g11:seq:sequence}{数列}为\omterm{def:g11:seq:arithmetic}{等差数列}。否则计算 $\frac{u_{n+1}}{u_n}$（各项
非零）并化简：若为常数 $q$，则为\omterm{def:g11:seq:geometric}{等比数列}。若两者都不是
常数，则该\omterm{def:g11:seq:sequence}{数列}既非\omterm{def:g11:seq:arithmetic}{等差数列}，也非\omterm{def:g11:seq:geometric}{等比数列}——切勿仅凭
前几项便下结论。
\end{method}

\begin{example}
对 $u_n = 3 \times 5^n$：
$\frac{u_{n+1}}{u_n} = \frac{3 \times 5^{n+1}}{3 \times 5^n} = 5$ 对所有
$n$，故为\omterm{def:g11:seq:geometric}{公比}为 $5$ 的\omterm{def:g11:seq:geometric}{等比数列}。对 $u_n = n^2$：
$u_1 - u_0 = 1$，但 $u_2 - u_1 = 3$，且
$\frac{u_1}{u_0}$ 甚至没有定义——既不是\omterm{def:g11:seq:arithmetic}{等差数列}，
也不是\omterm{def:g11:seq:geometric}{等比数列}。
\end{example}

\section{单调性}

\begin{definition}[单调数列]\label{def:g11:seq:monotone}
称\omterm{def:g11:seq:sequence}{数列} $(u_n)$ 为\emph{\omterm{def:g11:func:monotone}{递增}}，若 $u_{n+1} \geq u_n$ 对所有 $n$ 成立，
又称为\emph{\omterm{def:g10:functions:variations}{递减}}，若 $u_{n+1} \leq u_n$ 对所有 $n$ 成立。
\end{definition}

\begin{method}[研究单调性]\label{met:g11:seq:monotonicity}
研究 $u_{n+1} - u_n$ 的符号。对于各项为正的\omterm{def:g11:seq:sequence}{数列}，
也可改为比较 $\frac{u_{n+1}}{u_n}$ 与 $1$ 的大小。
\end{method}

\begin{example}
\omterm{def:g11:seq:arithmetic}{等差数列}在 $d \geq 0$ 时\omterm{def:g11:func:monotone}{递增}
（$u_{n+1} - u_n = d$），在 $d \leq 0$ 时\omterm{def:g10:functions:variations}{递减}。\omterm{def:g11:seq:geometric}{等比数列}
在 $u_0 > 0$ 且 $q > 1$ 时\omterm{def:g11:func:monotone}{递增}：
$u_{n+1} - u_n = u_0 q^n (q - 1) > 0$；当 $u_0 > 0$ 且 $0 < q < 1$ 时，它
\omterm{def:g10:functions:variations}{递减}。
\end{example}

\section{长期行为，直观描述}

$u_n$ 在 $n$ 变得很大时会怎样变化？对于\omterm{def:g11:seq:arithmetic}{等差数列}，
\omterm{def:g11:seq:arithmetic}{公差} $d > 0$ 时，各项 $u_0 + nd$ 最终会超过任何给定的数。
对于\omterm{def:g11:seq:geometric}{等比数列}，当 $0 < q < 1$ 时，各项
$u_0 q^n$ 趋向于 $0$：例如反复乘以 $0.9$，
任何初值都会被逐步削弱。而当 $q > 1$ 时，各项急剧增大，如
\cref{ex:g11:seq:chessboard}所示。

\begin{remark}
这些说法可以做到完全精确——“各项最终
会落在距 $0$ 任意给定距离之内”——并可加以证明。这正是
\emph{极限}理论，也是\cref{ch:g12:seq}的开篇主题。
\end{remark}

\section{练习}

\begin{exercise}[$\star$]\label{exo:g11:seq:1}
对每个\omterm{def:g11:seq:sequence}{数列}，计算 $u_1$、$u_2$、$u_3$：
\[
u_n = \frac{n}{n+1}; \qquad
u_0 = 5,\ u_{n+1} = 3u_n - 2; \qquad
u_n = (-1)^n\,n .
\]
\end{exercise}

\begin{exercise}[$\star$]\label{exo:g11:seq:2}
设 $(u_n)$ 为\omterm{def:g11:seq:arithmetic}{等差数列}，$u_0 = 7$，$d = -3$。计算 $u_{10}$ 与
$u_{25}$。设 $(v_n)$ 为\omterm{def:g11:seq:arithmetic}{等差数列}，$v_3 = 11$，$v_8 = 26$。求
\omterm{def:g11:seq:arithmetic}{公差}及 $v_0$。
\end{exercise}

\begin{exercise}[$\star$]\label{exo:g11:seq:3}
设 $(u_n)$ 为\omterm{def:g11:seq:geometric}{等比数列}，$u_0 = 5$，$q = 2$。计算 $u_8$。设 $(v_n)$ 为
正项\omterm{def:g11:seq:geometric}{等比数列}，$v_2 = 12$，$v_4 = 48$。求\omterm{def:g11:seq:geometric}{公比}
及 $v_0$。
\end{exercise}

\begin{exercise}[$\star$]\label{exo:g11:seq:4}
计算
\[
1 + 2 + 3 + \dots + 500, \qquad
4 + 7 + 10 + \dots + 61, \qquad
1 + \frac12 + \frac14 + \dots + \frac{1}{2^{10}} .
\]
\end{exercise}

\begin{exercise}[$\star$]\label{exo:g11:seq:5}
判断下列每个\omterm{def:g11:seq:sequence}{数列}是\omterm{def:g11:seq:arithmetic}{等差数列}、\omterm{def:g11:seq:geometric}{等比数列}，还是两者都不是：
\[
u_n = 4n - 1; \qquad
v_n = \frac{2^n}{3^{n+1}}; \qquad
w_n = n^2 + n .
\]
\end{exercise}

\begin{exercise}[$\star\star$]\label{exo:g11:seq:6}
某剧院共有 $20$ 排，第一排有 $16$ 个座位，每排座位数
比前一排多 $2$ 个座位。最后一排有多少
个座位？全剧院一共有多少个座位？
\end{exercise}

\begin{exercise}[$\star\star$]\label{exo:g11:seq:7}
某种细菌群每小时数量翻一番，正午时有 $500$
个细菌。晚上 8 点时有多少个？经过多少
个整小时后，数量首次超过一百万？（可依次尝试
$2$ 的各次幂求解。）
\end{exercise}

\begin{exercise}[$\star\star$]\label{exo:g11:seq:8}
每月，储户向一个账户存入 $100$ 欧元，该账户按
现有余额每月计付 $0.2\%$ 的利息（利息于存款前计入）。
记 $c_n$ 为第 $n$ 次存款后的
余额，于是 $c_1 = 100$，$c_{n+1} = 1.002\,c_n + 100$。计算 $c_2$
与 $c_3$，并说明 $(c_n)$ 为何既不是\omterm{def:g11:seq:arithmetic}{等差数列}，也不是\omterm{def:g11:seq:geometric}{等比数列}。
\end{exercise}

\begin{exercise}[$\star\star$]\label{exo:g11:seq:9}
研究下列\omterm{def:g11:seq:sequence}{数列}的单调性：
\[
u_n = n^2 - 8n \ (n \geq 0), \qquad
v_n = \frac{3^n}{n!}\ (n \geq 1),
\]
其中 $n! = 1 \times 2 \times \dots \times n$。（对于 $(v_n)$，比较
$\frac{v_{n+1}}{v_n}$ 与 $1$ 的大小。）
\end{exercise}

\begin{exercise}[$\star\star$]\label{exo:g11:seq:10}
一个\omterm{def:g11:seq:arithmetic}{等差数列}的前 $n$ 项和，其中 $u_0 = 3$，
$d = 4$，其和等于 $903$。求 $n$。（列出关于 $n$ 的
二次\omterm{def:g10:algebra:equation}{方程}，并用\cref{ch:g11:quad}中的方法。）
\end{exercise}

\begin{exercise}[$\star\star\star$]\label{exo:g11:seq:11}
设 $u_0 = 3$，$u_{n+1} = 2u_n - 1$。
\begin{enumerate}
  \item 计算 $u_1, u_2, u_3$，并猜想 $u_n$ 的通项公式。
  \item 令 $v_n = u_n - 1$。证明 $(v_n)$ 为\omterm{def:g11:seq:geometric}{等比数列}，并求其
        \omterm{def:g11:seq:geometric}{公比}与首项。
  \item 由此推出 $u_n$ 的通项公式，并验证你的猜想。
\end{enumerate}
\end{exercise}

\section{问题：梵天塔与斐波那契的兔子}

\begin{problem}[{周末问题——两个传奇递推式：
令世界终结的塔，像黄金一样增长的数列，
以及驯服贷款的辅助数列技巧}]
\label{pb:g11:seq:1}
数学的民间传说中，有两个\omterm{def:g11:seq:sequence}{数列}独领风骚。一个计算
梵天塔的移动次数——六十四个金盘的
搬运，据传说将终结世界。另一个\omterm{def:g11:seq:sequence}{数列}计算
斐波那契的兔子，并隐藏着黄金分割比。两者
皆非\omterm{def:g11:seq:arithmetic}{等差数列}，亦非\omterm{def:g11:seq:geometric}{等比数列}——二者都会屈服于
本章的武器：递推式、\omterm{def:g11:seq:geometric}{等比数列}求和
（\cref{thm:g11:seq:geomsum}）以及
辅助\omterm{def:g11:seq:sequence}{数列}技巧（见\cref{exo:g11:seq:11}），它还能算出你的房贷。

\medskip
\noindent\textbf{第一部分——梵天塔。}
谜题：$n$ 个大小\omterm{def:g10:functions:variations}{递减}的圆盘叠放在柱 A 上；
每次移动一个圆盘，将整叠圆盘移到柱 C 上，
且不得把较大的圆盘放在较小的圆盘之上（柱 B 可以辅助）。记 $h_n$ 为
所需的最少移动次数。
\begin{enumerate}
  \item 用硬币动手试一试，并记录 $h_1$、$h_2$、$h_3$。
  \item 说明递推式背后的思路：
        $h_{n+1} = 2h_n + 1$ 中，最大圆盘移动前后
        分别必须发生什么？
  \item 借助
        \cref{exo:g11:seq:11}中的技巧求解递推式：令 $v_n = h_n + 1$，证明
        $(v_n)$ 为\omterm{def:g11:seq:geometric}{等比数列}，并由此得出 $h_n = 2^n - 1$。
  \item 传说中的塔有 $64$ 个圆盘，僧侣们每秒
        移动一个圆盘。利用 $2^{10} = 1024 \approx 10^3$，
        估算以年计的搬运时间（一年约为
        $3 \times 10^7$ 秒；可参照
        \cref{ex:g11:seq:chessboard}，同一个庞然大物换了一个
        故事出现）。是否需要担心？
  \item 为何任何策略都无法少于 $2^n - 1$ 次移动？论证
        \emph{任一}移动方案都满足
        $h_{n+1} \geq 2 h_n + 1$：顶部的
        $n$ 个圆盘在最底层圆盘移动之前与之后
        各须满足什么？
\end{enumerate}

\medskip
\noindent\textbf{第二部分——斐波那契\omterm{def:g11:seq:sequence}{数列}。}
定义 $F_1 = F_2 = 1$，且 $F_{n+2} = F_{n+1} + F_n$（每一项都是前两项
之和——即初中分册中的“数节拍”规则，
只是现在有了它的欧洲名字）。
\begin{enumerate}[resume]
  \item 列出 $F_1$ 至 $F_{12}$。
  \item 证明 $(F_n)$ 既不是\omterm{def:g11:seq:arithmetic}{等差数列}，也不是\omterm{def:g11:seq:geometric}{等比数列}，
        但从 $n = 2$ 起它严格\omterm{def:g11:func:monotone}{递增}
        （可借助\cref{met:g11:seq:monotonicity}与递推关系）。
  \item 证明求和恒等式
        \[
        F_1 + F_2 + \dots + F_n = F_{n+2} - 1
        \]
        用裂项相消法：将每个 $F_k$ 写成
        $F_{k+2} - F_{k+1}$，观察求和式如何消去。验证
        它在 $n = 6$ 时成立。
  \item 证明平方和恒等式
        $F_1^2 + F_2^2 + \dots + F_n^2 = F_n F_{n+1}$，
        用裂项相消法，其中
        $F_k F_{k+1} - F_{k-1} F_k = F_k^2$。验证
        $n = 4$ 时成立。（图示：边长为
        $1, 1, 2, 3, 5, \dots$ 的正方形可以拼成一个矩形——即著名的
        斐波那契螺旋的骨架。）
  \item 卡西尼恒等式为
        $F_{n+1} F_{n-1} - F_n^2 = (-1)^n$。验证它在
        $n = 4, 5, 6$ 时成立，并认出
        初中分册面积题中“消失正方形”戏法的内在动力，
        即该题所用的那一招。
  \item 由递推关系证明
        $F_{n+2} \geq 2 F_n$：斐波那契\omterm{def:g11:seq:sequence}{数列}每两步至少翻一番，
        即它至少与一个\omterm{def:g11:seq:geometric}{公比}为 $\sqrt2$ 的
        \omterm{def:g11:seq:geometric}{等比数列}增长得一样快。
  \item 计算比值 $r_n = \frac{F_{n+1}}{F_n}$ 在
        $n = 3$ 至 $10$ 的值（保留三位小数）。承认它们
        趋于某极限 $L$，将关系式
        $r_{n+1} = 1 + \frac{1}{r_n}$ 取极限并求解：
        兔子所崇拜的，是\cref{pb:g10:algebra:1}
        中的哪一个数？
\end{enumerate}

\medskip
\noindent\textbf{第三部分——辅助\omterm{def:g11:seq:sequence}{数列}技巧，
在银行里。}
\begin{enumerate}[resume]
  \item 推广\cref{exo:g11:seq:11}：当
        $u_{n+1} = a\,u_n + b$，$a \neq 1$ 时，令
        $\ell = \frac{b}{1 - a}$（\omterm{pb:g10:functions:1}{不动点}）。证明
        $v_n = u_n - \ell$ 是\omterm{def:g11:seq:geometric}{公比}为 $a$ 的\omterm{def:g11:seq:geometric}{等比数列}，并
        由此得出
        $u_n = a^n (u_0 - \ell) + \ell$。
  \item 一笔贷款：$10\,000$ 欧元，月利率 $1\,\%$，
        每月偿还 $300$ 欧元，因此债务满足
        $d_{n+1} = 1.01\,d_n - 300$。运用第 13 题（先求
        \omterm{pb:g10:functions:1}{不动点}！），得出 $d_n$ 的显式表达式。
  \item 借助计算器，求出债务首次被还清的月份，
        以及累计偿还的总额。借款本身
        一共花费了多少？
  \item 一座有 $50\,000$ 居民的小镇每年增长 $2\,\%$，
        此外每年还迎来 $1\,000$ 名新居民：
        $p_{n+1} = 1.02\,p_n + 1000$。给出显式表达式
        并求 $10$ 年后的人口数。
\end{enumerate}

\medskip
\noindent\textbf{第四部分——两大王族。}
\begin{enumerate}[resume]
  \item 计算 $1 + 2 + 3 + \dots + 1000$
        （\cref{thm:g11:seq:intsum}——初中分册中的“小高斯”求和，
        现在正式成立），以及
        $1 + 2 + 4 + \dots + 2^{19}$
        （\cref{thm:g11:seq:geomsum}）。
  \item 计算\omterm{def:g11:seq:arithmetic}{等差数列}的和：
        $7, 12, 17, \dots, 502$（共有多少项？）。
  \item 储蓄计划：每月存入 $100$ 欧元，
        按月利率 $0.5\,\%$ 计息；第 $n$ 次存款后，余额为
        $100\left(1.005^{n-1} + \dots + 1.005 + 1\right)$。
        求 $5$ 年后的余额（$n = 60$）。
  \item 总结——\omterm{def:g11:seq:sequence}{数列}驯服者的工具箱：显式描述与递推描述的对照；
        两大王族及其求和公式；
        把仿射递推化为\omterm{def:g11:seq:geometric}{等比数列}的辅助\omterm{def:g11:seq:sequence}{数列}；
        以及斐波那契\omterm{def:g11:seq:sequence}{数列}——它是第一位不属于两大王族的“公民”，
        目前借恒等式将其驯服，
        有待矩阵（十二年级）与极限来完整捕获。
        每项各用一句话概括。
\end{enumerate}
\end{problem}
